\chapter{Impacto Social e Econômico dos Ecossistemas Cognitivos}

\begin{quote}
\textit{``A pergunta não é se a IA vai mudar o mundo, mas quem controla essa mudança e para o benefício de quem.''} --- Kate Crawford, Atlas of AI (2021)
\end{quote}

\section{A Economia da Atenção e o Capitalismo de Vigilância}

Shoshana Zuboff (2019) \cite{zuboff2019age} denominou \textbf{capitalismo de vigilância} o sistema econômico que reivindica a experiência humana como matéria-prima gratuita para extração, predição e vendas. Tim Wu (2016) \cite{wu2016attention} documentou como a \textbf{economia da atenção} transformou publicidade, mídia e política, onde a atenção humana é um recurso finito e não-renovável disputado ferozmente por plataformas que otimizam engajamento, não bem-estar.

O valor de uma plataforma pode ser modelado como:

\begin{equation}
\text{Valor}_{\text{plataforma}} = \sum_{u \in \text{usuários}} \text{Tempo}(u) \times \text{CPM}(u) \times \text{Precisão}_{\text{targeting}}(u)
\label{eq:platform_value}
\end{equation}

A resposta regulatória inclui GDPR (UE, 2018), LGPD (Brasil, 2020) e o AI Act (UE, 2024). O OpenCode incorpora conformidade via Agente 32 (Ética \& Open Science), Agente 34 (Conflitos \& Similaridade), e processamento local via Ollama para dados sensíveis.

\subsection{PRÁXIS 2.1 — Auditor de Extração de Dados}

\begin{lstlisting}[language=Python, caption={SurveillanceAuditor — verificador de conformidade LGPD/GDPR}]
class SurveillanceAuditor:
    """Audita o fluxo de dados em ecossistemas de agentes."""
    def __init__(self, agent_id):
        self.agent_id = agent_id
        self.data_flows = []
        self.alerts = []
    
    def log_extraction(self, source, data_type, volume_bytes, purpose):
        self.data_flows.append({
            'timestamp': datetime.utcnow().isoformat(),
            'source': source, 'data_type': data_type,
            'volume_bytes': volume_bytes, 'purpose': purpose
        })
    
    def audit_privacy(self, consent_required=True):
        violations = []
        for flow in self.data_flows:
            if flow['purpose'] not in ['essential', 'improvement', 'consented']:
                violations.append({'violation': 'Propósito não declarado'})
            if flow['volume_bytes'] > 1_000_000 and flow['data_type'] == 'personal':
                self.alerts.append(f"ALERTA: Coleta massiva ({flow['volume_bytes']} bytes)")
        return {'compliant': len(violations) == 0, 'violations': violations}
\end{lstlisting}

\section{Democratização da IA via Ecossistemas Abertos}

O movimento open source, articulado por Stallman (1985) \cite{stallman2002free} e Raymond (1999) \cite{raymond1999cathedral}, estabeleceu os princípios que guiam a democratização da IA. A partir de 2023, modelos abertos como LLaMA (Touvron et al., 2023) \cite{touvron2023llama}, Mistral (Jiang et al., 2023) \cite{jiang2023mistral} e DeepSeek \cite{deepseek2024v2} desafiaram o domínio das Big Tech. A plataforma Hugging Face (Wolf et al., 2020) \cite{wolf2020transformers} tornou-se o ``GitHub da IA'' com mais de 500.000 modelos.

O OpenCode estende este movimento: ecossistema aberto, auditável e extensível para produção científica Qualis A1.

\subsection{PRÁXIS 2.2 — Fine-Tuning Democrático com QLoRA}

\begin{lstlisting}[language=Python, caption={QLoRA: fine-tuning de LLMs em GPU de consumo}]
from transformers import AutoModelForCausalLM, BitsAndBytesConfig
from peft import LoraConfig, get_peft_model, prepare_model_for_kbit_training

def democratize_finetuning(base_model="mistralai/Mistral-7B-v0.1"):
    bnb_config = BitsAndBytesConfig(
        load_in_4bit=True, bnb_4bit_quant_type="nf4",
        bnb_4bit_compute_dtype=torch.bfloat16, bnb_4bit_use_double_quant=True
    )
    model = AutoModelForCausalLM.from_pretrained(base_model, quantization_config=bnb_config, device_map="auto")
    model = prepare_model_for_kbit_training(model)
    lora_config = LoraConfig(r=16, lora_alpha=32, target_modules=["q_proj","k_proj","v_proj","o_proj"])
    model = get_peft_model(model, lora_config)
    # trainable params: ~4.2M de 7.2B (0.058%) — roda em RTX 3060 12GB
    return model
\end{lstlisting}

\section{Impacto no Mercado de Trabalho e Educação}

Frey e Osborne (2017) \cite{frey2017future} estimaram 47\% dos empregos nos EUA em risco de automação. Acemoglu e Restrepo (2022) \cite{acemoglu2022tasks} refinaram com modelo de tarefas: a automação redistribui tarefas — o impacto depende se novas tarefas são criadas mais rápido que as antigas automatizadas.

Na educação, o ``2 Sigma Problem'' de Bloom (1984) \cite{bloom19842sigma} — tutoria individual produz 2 desvios-padrão de melhoria — encontra nos tutores de IA sua realização técnica. O pipeline MASWOS do OpenCode (Capítulo 11) democratiza a produção acadêmica de alto rigor.

\section{Token Economy como Novo Paradigma}

Voshmgir (2020) \cite{voshmgir2020token} propõe tokens como sinais multidimensionais: acesso, voto, propriedade e reputação. No OpenCode (Capítulo 10): staking proporcional à confiança, recompensa por sucesso (gate SDD + BehavioralGate), slashing por falha, e fee market para priorização.

\section{Ética e Sustentabilidade}

Floridi et al. (2021) \cite{floridi2021ai4people} propôs o framework AI4People com 5 princípios: Beneficência, Não-maleficência, Autonomia, Justiça e Explicabilidade — todos operacionalizados no OpenCode via BehavioralGate, provenance tracking, e Token Economy.

O custo ambiental da IA (GPT-3: 552 t CO$_2$eq; Patterson et al., 2021 \cite{patterson2021carbon}) é mitigado pelo OpenCode via Ollama local (inferência em hardware de consumo), quantização (4-bit, redução de 75\%+) e reuso metacognitivo (memória compartilhada evita recomputação).

O ecossistema metacognitivo não é apenas arquitetura de software — é uma proposta de \textbf{infraestrutura social} para a era da IA.

\begin{thebibliography}{99}
\bibitem{zuboff2019age} Zuboff, S. (2019). \textit{The Age of Surveillance Capitalism}. PublicAffairs.
\bibitem{wu2016attention} Wu, T. (2016). \textit{The Attention Merchants}. Knopf.
\bibitem{stallman2002free} Stallman, R. (2002). \textit{Free Software, Free Society}. GNU Press.
\bibitem{raymond1999cathedral} Raymond, E. S. (1999). \textit{The Cathedral and the Bazaar}. O'Reilly Media.
\bibitem{wolf2020transformers} Wolf, T., et al. (2020). Transformers: State-of-the-Art NLP. \textit{EMNLP 2020}. \doiurl{10.18653/v1/2020.emnlp-demos.6}
\bibitem{frey2017future} Frey, C. B., \& Osborne, M. A. (2017). The Future of Employment. \textit{Technological Forecasting and Social Change}, 114, 254--280. \doiurl{10.1016/j.techfore.2016.08.019}
\bibitem{acemoglu2022tasks} Acemoglu, D., \& Restrepo, P. (2022). Tasks, Automation, and U.S. Wage Inequality. \textit{Econometrica}, 90(5), 1973--2016. \doiurl{10.3982/ECTA19815}
\bibitem{voshmgir2020token} Voshmgir, S. (2020). \textit{Token Economy: How the Web3 Reinvents the Internet}. Token Kitchen.
\bibitem{floridi2021ai4people} Floridi, L., et al. (2021). How to Design AI for Social Good. \textit{Science and Engineering Ethics}, 27(2). \doiurl{10.1007/s11948-021-00295-1}
\bibitem{patterson2021carbon} Patterson, D., et al. (2021). Carbon Emissions and Large Neural Network Training. \textit{arXiv:2104.10350}. \doiurl{10.48550/arXiv.2104.10350}
\bibitem{bloom19842sigma} Bloom, B. S. (1984). The 2 Sigma Problem. \textit{Educational Researcher}, 13(6), 4--16. \doiurl{10.3102/0013189X013006004}
\bibitem{crawford2021atlas} Crawford, K. (2021). \textit{Atlas of AI}. Yale University Press.
\bibitem{touvron2023llama} Touvron, H., et al. (2023). Llama 2: Open Foundation Models. \textit{arXiv:2307.09288}. \doiurl{10.48550/arXiv.2307.09288}
\bibitem{jiang2023mistral} Jiang, A. Q., et al. (2023). Mistral 7B. \textit{arXiv:2310.06825}. \doiurl{10.48550/arXiv.2310.06825}
\bibitem{deepseek2024v2} DeepSeek-AI. (2024). DeepSeek-V2. \textit{arXiv:2405.04434}. \doiurl{10.48550/arXiv.2405.04434}
\end{thebibliography}
